\documentclass{jsarticle}
\usepackage{amsmath}
\usepackage{amssymb}
\begin{document}
\title{}
\author{}
\date
\maketitle

坂本龍一は私にとって、最高の恩師だ。
まず言うには、スコラの音楽学校での、曲の作り方を直感を通して、
私よりは、上の力量の人たち、私が高望みをしているレベルの人たち
を出演させて、中にあやさんの名前に近い、神社関係の名前の女性
が、名前は善光寺彩子さんで、坂本龍一に近い曲を作っていたのを、
このときに各人の曲を坂本龍一が、その作曲している曲の使っている
音階を聞き始めただけで一瞬にして言い当てているこの才能を、
石川先生がバッハのロ短調のミサのことを
私が言われているのを、退院したあとに、スガナミ楽器のバイオリンが
専門の女性に音階の仕組み自体の機知を教えてもらっていたのを、
最近になって、ファインダーに書き書きしていたら、
ロ短調はB♭ 、ラストエンペラーのE♭はFがない、演歌のG♭はFがない、
鋭い共鳴はF♯はGがない、このB♭とE♭とG♭はきれいすぎて、東風はB♭、
なぜ坂本龍一の曲は、不協和音がものすごく多いのにすごくきれいかが分かった。
この音階を知る前から、sweet revengeをヒントに直感で自動書記みたいな
能力を順子お姉さんと睦世お姉さんに教えてもらった、ハノンを猛練習した
能力がトリガーとなって、音階を知らないために、音階がＨからＺへと
音楽が音学の範囲を超えて、曲がＣをベースにしていながら、不協和音みたいな
曲なのに和音のオンパレードになっていて、自分でも気に入っている。
坂本龍一には本当に１００歳を超えてずっと私の好きな曲とわけのわからない、
これが音楽の研究とおもうが、現代音楽を超えて、ずっと１００歳を超えて、
曲をずっと作っていてほしい。
ラストエンペラーのE♭は、皇帝の傀儡のもの虚しいが、離婚のシーンと気品と
仕方なく妥協するが、屈服する哀れなのに、全体がきれいな分散和音と
不協和音が流れている。自分でも弾いてみたいが、楽譜を見ると、
弾けるか！と思っていたが、音階の秘密を知るとなるほどと言えた。
Forbidden cloursは、デヴィッド・シルビアンのダンディーな声もいいが、
順子お姉さんが、楽譜を速読しながらのピアンで聞いていると、
戦場のメリークリスマスより、こっちのほうがものすごくきれい。
姉が久石譲さんのVenusを弾いていると、本人よりも言うのがやりきれないが、
ものすごくきれい。実働先生が順子お姉さんにこのピアノにベタぼれしていたらしく、
数学と国語が順子お姉さんは１０段階中１０だったらしく、美鈴ケ丘中学校３年７
組のクラスは、ここは天国ではないかとおもうぐらい、友達は話しかけてこなかったが、
ものすごく幸せだった。順子お姉さんは中学校のときに、バレーボールが
出来て、勉強が出来て、ピアノまで出来て、河合塾はエンリッチクラスで、
美術が出来て、水泳が出来て、姿が、今は写真写りが悪いが、洋服のファッションセンス
が抜群で、性格だけが欠点であるのは、姉たちが凄すぎたので、
私の取り柄は剣道だけだったらしい。姉たち二人が私の剣道をしているのをみて、
自分たちもやりたいと言っていたが、私は薦められずに、やらないほうがいいよと
言っていた。福山市に来て、美鈴ケ丘での習い事を全部やりなおした。
全部出来た。数学と音楽の能力は相関関係があることをＺ会の英語で知った。
モーツァルトがそうだったらしい。
東風は、Ｂ♭でありながら、ＣとＥ♭で共鳴できているのに不協和音であり、
いけいけな曲でありながら、聴くのを我慢しなくてはならずにいて、
それでいて、きれいな連弾でもあり、Happy endとは違う曲だが、
速読の能力で作られている。砂漠とオアシスの狭間と吹き抜ける風が
共感覚と誰でも体験できる曲になっている。
\end{document}
